\section{Discussion and Limitations}
\label{sec:discussion}

\upd{We discuss the core insights, metric divergences, architectural principles,
and limitations established by our experiments.}

\paragraph{Privileged parameter visibility and modality competition.}
\upd{A central design question in multimodal architectures trained alongside
generative models is where privileged structural information should enter
the network.} \upd{The procedural parameter stream $\theta$ is the only modality
that states target phenotypic traits verbatim, yet it is never available at
deployment in the field.} \upd{Across our controlled ablations
(Sec.~\ref{sec:exp-params}), exposing parameter tokens to the encoder
during pretraining degraded the quality of learned sensor representations
in both sorghum and maize.} \upd{In sorghum, every model arm that saw parameter
tokens in the encoder learned inferior visual features, with the
degradation on the size traits growing with the stream's visibility
(Tab.~\ref{tab:token_budget}).}

\upd{This behavior can be understood through the lenses of shortcut learning,
modality competition~\citep{peng2022ogm}, and negative gradient
transfer~\citep{yu2020pcgrad}.} \upd{When ground-truth parameter tokens are
visible to the transformer encoder, they directly encode target geometry.} \upd{A
plausible account is that the encoder learns a low-cost shortcut through
these parameter tokens, starving the visual (RGB, depth) and point cloud
streams of gradient updates; when parameter tokens are then absent at
inference, the sensor representations prove impoverished.}

\upd{The species-specific control arms (Tab.~\ref{tab:token_budget}) reveal an
additional nuance:}
\begin{itemize}
\item \upd{ \textbf{Sorghum}: Keeping parameter tokens strictly as a decoder
  target ($\rho_\theta = 1.0$) continues to lower sensor latent quality
  ($R^2$ drops from $0.977$ to $0.966$ on height).} \upd{Sorghum morphological
  parameters are tightly coupled (effective rank two); a plausible
  mechanism is gradient competition during multi-task decoding that pulls
  sensor representations toward global plant size, the leading principal
  component of the sorghum latent.} \item \upd{ \textbf{Maize}: Keeping parameters strictly as a decoder target
  ($\rho_\theta = 1.0$) or restricting visibility to 20\% improves sensor
  latent probing across all eleven traits ($R^2 = 0.776$ vs.\ $0.633$ on
  height).} \upd{Maize parameters contain 11 independent morphological
  dimensions; supervising them as a generation target provides structured
  multi-task decoder gradients, consistent with a regularizing effect on
  sensor latents without encoder shortcuts.}
\end{itemize}

\upd{These findings yield an architectural design rule for multimodal masked
autoencoders trained on synthetic or procedural data: \emph{privileged
generative parameters should never be largely visible to the encoder during
pretraining}.} \upd{Setting parameter tokens strictly as a decoder reconstruction
target ($\rho_\theta = 1.0$) or applying heavy gating
($\rho_\theta \le 0.20$) prevents encoder shortcut learning while
preserving generative parameter output capabilities; whether the parameters
help even as a pure reconstruction target depends on the generator.}

\paragraph{Metric divergence on slender geometry.}
\upd{A critical empirical finding of this work is the structural divergence
between nearest-neighbor point cloud metrics (such as Chamfer distance) and
downstream phenotypic trait probing on slender geometry
(Sec.~\ref{sec:exp-metrics}, Fig.~\ref{fig:metric_divergence}).} \upd{Standard
Chamfer distance evaluates surface proximity by computing average
nearest-neighbor point displacements:}
\begin{equation}
  \mathrm{CD}(P, \hat P) = \frac{1}{|P|}\sum_{p\in P}\min_{q\in\hat P}\|p-q\|_2^2 + \frac{1}{|\hat P|}\sum_{q\in\hat P}\min_{p\in P}\|p-q\|_2^2.
  \label{eq:chamfer_def}
\end{equation}
\upd{Because crop organs (such as maize and sorghum leaves) are thin,
high-aspect-ratio surfaces, a predicted point cloud can achieve low Chamfer
distance simply by clustering points onto bulky central stalks while
omitting entire leaf blades.} \upd{In contrast, downstream linear probes fitted
on frozen encoder embeddings directly measure whether latent
representations retain the structural parameter information required for
phenotyping.}

\upd{Throughout our experiments, Chamfer distance and linear probes rank model
variants differently:}
\begin{enumerate}
\item \upd{ \textbf{Depth stream addition}: Adding metric depth to RGB
  photographs reduces Chamfer distance by 28\% in sorghum and 25\% in
  maize, but yields negligible change in downstream probe $R^2$ scores
  (Tab.~\ref{tab:modality_value_add}).} \upd{Depth ordering provides spatial
  surface proximity for decoders without expanding high-level trait
  information in the encoder.} \item \upd{ \textbf{Data scaling}: Increasing pretraining dataset volume from
  1{,}000 to 10{,}000 plants reduces sorghum Chamfer error by 21\%
  ($3.13 \to 2.48 \times 10^{-3}$), whereas sorghum probe scores saturate
  by 10{,}000 plants (Tab.~\ref{tab:scaling_ablations}).} \item \upd{ \textbf{Partial camera visibility}: Restricting point cloud inputs
  to a single camera perspective increases Chamfer error $2.6\times$ to
  $6.5\times$, whereas refitted linear probes of the fused arms lose at
  most $0.030$ in height $R^2$ (Tab.~\ref{tab:views}).}
\end{enumerate}

\upd{These disagreements demonstrate that Chamfer distance alone is an
insufficient metric for evaluating slender 3D shape reconstruction.} \upd{We
advocate evaluating structural downstream trait probes alongside Chamfer
distance, and incorporating distribution-matching losses (such as Sinkhorn
optimal transport divergence, Eq.~\eqref{eq:point-loss-sinkhorn}) during
training to force decoders to cover full organ geometry rather than
collapsing points onto central stalks.}

\paragraph{Partial coverage, camera viewpoint, and row occlusion.}
\upd{Sensor observations in agricultural field campaigns are inherently partial.} \upd{Our experiments analyze three aspects of partial sensing:}
\begin{itemize}
\item \upd{ \textbf{Viewpoint elevation dominance}: Across isolated camera
  perspectives, camera elevation angle dominates single-image
  reconstruction difficulty, accounting for $\eta^2 = 0.24\text{--}0.31$ of
  per-plant reconstruction variance (Fig.~\ref{fig:views}).} \upd{Steep top-down
  and bottom-up perspectives suffer from severe self-occlusion where leaves
  stack vertically along the optical axis, whereas side-on views expose
  full blade silhouettes.} \item \upd{ \textbf{Covariate shift under partial visibility}: While phenotypic
  trait information persists in encoder representations under partial
  surface coverage (refitted 1-camera height $R^2 = 0.960$), directly
  transferring linear probes trained on 100\% complete clouds to partial
  inputs results in performance failure ($R^2 \to -0.26$,
  Fig.~\ref{fig:partial}).} \upd{Partial visibility induces a systematic
  covariate shift in latent feature space.} \upd{Consequently, downstream
  phenotyping readouts deployed in the field must be refit or calibrated on
  sensor-matched partial coverage inputs.} \item \upd{ \textbf{Implicit row occlusion reasoning}: In planted rows,
  neighboring foliage accounts for up to two-thirds of input point clouds
  and pixel areas.} \upd{Standard models trained on isolated plants suffer
  catastrophic failure under row occlusion because decoders attempt to
  incorporate neighboring leaves into target plant reconstructions
  (Tab.~\ref{tab:regime_3}).} \upd{Fine-tuning decoders to predict target main
  plant geometry behind occluding neighbors removes 92\% of
  neighbor-induced occlusion error in maize and 85\% in sorghum, forcing
  the network to perform implicit occlusion reasoning and segment target
  organs from inter-plant contact foliage.}
\end{itemize}

\upd{Importantly, procedural parameter outputs ($\hat\theta$) demonstrate
superior robustness to partial visibility compared to raw point cloud
completions.} \upd{Because the procedural generator $G(\hat\theta)$ acts as an
explicit biological prior, re-rendered procedural meshes generate complete
3D geometry even when input sensor observations capture only a fraction of
the plant surface.}

\paragraph{Capabilities and boundaries of cross-modal distillation.}
\upd{Cross-modal teacher--student distillation (Sec.~\ref{sec:method-distill})
converts a multimodal pretraining artifact into a fast, deployable
single-sensor inference model.} \upd{Distilling a full-modality teacher
($\mathrm{RGB}+\mathrm{D}+\mathrm{PC}+\theta$) into single-sensor student
networks reduces overall cross-modal generation loss by 30.2\% in sorghum
and 43.5\% in maize (Tab.~\ref{tab:anysubset}).} \upd{Distillation gains are
concentrated on the weakest input modalities (photograph-to-cloud Chamfer
error falls by 61\%) and the most challenging plant geometries
(Fig.~\ref{fig:views}B).} \upd{Furthermore, supervising student models across all
four decoder targets simultaneously outperforms a specialized
parameter-only regression student \todo{(first release; no final-data
specialist was trained)}, consistent with auxiliary visual reconstruction
tasks acting as regularizing gradients for parameter regression.}

\upd{However, distillation exhibits clear boundary limits:}
\begin{enumerate}
\item \upd{ \textbf{Viewpoint robustness}: Distillation lowers absolute
  reconstruction error across all camera views but does not flatten the
  U-shaped elevation difficulty curve (Fig.~\ref{fig:views}A).} \upd{Viewpoint
  robustness cannot be acquired purely through distillation; it requires
  balanced elevation sampling during pretraining data capture
  (Fig.~\ref{fig:e5views}).} \item \upd{ \textbf{Unobservable parameter instability}: Parameterizations
  submitted to inverse regression must be restricted to traits that field
  sensors can physically resolve.} \upd{In early experimental iterations
  containing unobservable leaf waviness parameters, distillation caused
  these unconstrained parameter heads to develop severe view-dependent
  instability without improving trait skill (skill $\approx 0$).} \upd{Restricting procedural parameter representations to observable
  physiological attributes (insertion heights, blade lengths, widths,
  angles, and rolls) ensures stable cross-modal distillation.}
\end{enumerate}

\paragraph{Limitations and future directions.}
\upd{We acknowledge several limitations in our current study:}
\begin{itemize}
\item \upd{ \textbf{Single-seed ablation budgets}: Controlled ablation
  experiments (Tabs.~\ref{tab:modality_value_add}--\ref{tab:scaling_ablations})
  were executed under a standardized compute budget of 197{,}400 optimizer
  steps using a single seed per arm.} \upd{Run-to-run baseline variance estimated
  from near-duplicate runs is $\approx 2\%$ in Chamfer distance and
  $\le 0.011$ in probe $R^2$; the non-monotone maize scaling differences,
  though significant under the paired bootstrap, could still reflect
  run-to-run variance.} \item \upd{ \textbf{Procedural generator expressivity boundaries}: An inherent
  limitation of inverse procedural modeling is that recovered
  representations are bounded by the expressivity of the underlying
  generator $G$.} \upd{If a physical plant exhibits morphological deformations
  unrepresented in $G$ (such as severe disease lesions, structural
  breakage, or tassel droop, which the generator renders without learnable
  variation), the framework fits the nearest valid procedural parameter
  configuration.} \upd{Expanding procedural generators to support broader
  developmental malformations represents an important avenue for future
  research.} \item \upd{ \textbf{Synthetic renderer regularities}: Synthetic rendering
  pipelines can introduce unintended canonical cues (such as fixed leaf
  plane azimuth alignments) that pose-aware baselines may exploit as
  shortcuts.} \upd{Rigorous evaluation on real-world field captures
  (Sec.~\ref{sec:exp-real}) remains essential to verify true sim-to-real
  domain transfer.}
\end{itemize}

\upd{By releasing \dataset{} (comprising 30{,}000 synthetic plants across maize
and sorghum with paired multimodal captures, physically simulated row
occlusion, and exact procedural parameter ground truth), we aim to provide
the 3D vision and digital agriculture communities with a benchmark for
evaluating multimodal completion and inverse procedural modeling on
slender, heavily occluded geometry.}
